\documentclass[11pt]{article}

\usepackage{amsmath, amssymb}
\usepackage{geometry}
\geometry{margin=1in}

\title{
CSE400 -- Fundamentals of Probability in Computing\\
Lecture 11 \& 12: Transformation of Random Variables
}
\author{Instructor: Dhaval Patel, PhD}
\date{February 10, 2026}

\begin{document}

\maketitle

%------------------------------------------------

\section{Introduction and Objectives}

The primary objective of this lecture is the learning of transformation techniques for random variables.  
This involves determining the probability distribution of a new random variable $(Y)$ derived from an existing one $(X)$.

\subsection{Core Assumptions}

\begin{itemize}
    \item The Probability Density Function (PDF), $f_X(x)$, of the original random variable $X$ is known \textit{a priori}.
    \item The Cumulative Distribution Function (CDF), $F_X(x)$, of $X$ is also known.
    \item A new random variable $Y$ is defined by a functional relationship:
    \[
    y = g(x).
    \]
\end{itemize}

\subsection{Objectives}

\begin{itemize}
    \item Determine how to find the new distribution $f_Y(y)$ or $F_Y(y)$.
    \item Understand joint transformations for two random variables.
\end{itemize}

%------------------------------------------------

\section{Transformation of a Single Random Variable}

The derivation of the new PDF follows a systematic two-step process.

\subsection{Step 1: Find the Cumulative Distribution Function (CDF)}

The CDF of the new variable $Y$, $F_Y(y)$, is the probability that $Y$ is less than or equal to a value $y$:

\[
F_Y(y) = \Pr(Y \le y) = \Pr(g(x) \le y)
\]

If the function $g(x)$ is monotonically increasing:

\[
F_Y(y) = \Pr(x \le g^{-1}(y)) = F_X(g^{-1}(y))
\]

If the function $g(x)$ is decreasing:

\[
F_Y(y) = \Pr(x \ge g^{-1}(y)) = 1 - F_X(g^{-1}(y))
\]

%------------------------------------------------

\subsection{Step 2: Differentiate to Find the PDF}

To find the PDF $f_Y(y)$, differentiate the CDF with respect to $y$:

\[
f_Y(y) = \frac{d}{dy}\Big[F_X(g^{-1}(y))\Big]
\]

Using the chain rule, this results in:

\[
f_Y(y) = f_X(g^{-1}(y)) \cdot \frac{d}{dy}g^{-1}(y)
\]

Or more generally, accounting for both increasing and decreasing functions:

\[
f_Y(y) = f_X(x)\left|\frac{dx}{dy}\right|_{x=g^{-1}(y)}
\]

%------------------------------------------------

\section{Illustrative Example: Uniform to Sine Transformation}

\textbf{Problem:}  
Given $X$ is a uniformly distributed random variable on the interval $(-1,1)$, find the PDF of

\[
Y = \sin\left(\frac{\pi x}{2}\right).
\]

\subsection{Solution}

\textbf{Original PDF of $X$:}

\[
f_X(x) = \frac{1}{2}, \qquad -1 < x < 1,
\]
and $0$ otherwise.

\textbf{Inverse Function:}

\[
x = \frac{2}{\pi}\sin^{-1}(y).
\]

\textbf{Derivative:}

\[
\frac{dx}{dy} = \frac{2}{\pi}\cdot \frac{1}{\sqrt{1-y^2}}.
\]

\textbf{Resulting PDF:}

\[
f_Y(y) = \frac{1}{2}\cdot \frac{2}{\pi}\cdot \frac{1}{\sqrt{1-y^2}}
       = \frac{1}{\pi\sqrt{1-y^2}}.
\]

\textbf{Limits:}

For $x=-1$, $y=\sin(-\pi/2)=-1$.  
For $x=1$, $y=1$.  

Thus,

\[
-1 < y < 1.
\]

%------------------------------------------------

\section{Function of Two Random Variables}

This section addresses joint transformations where

\[
Z = g(X,Y).
\]

%------------------------------------------------

\subsection{Detailed Derivation: Sum of Two RVs ($Z=X+Y$)}

To find the PDF $f_Z(z)$ for the sum $Z=X+Y$:

\subsubsection{Determine the CDF}

\[
F_Z(z) = \Pr(X+Y \le z).
\]

\subsubsection{Set up the Double Integral}

\[
F_Z(z) = \int_{-\infty}^{\infty}\int_{-\infty}^{z-y}
f_{XY}(x,y)\,dx\,dy.
\]

%------------------------------------------------

\subsubsection{Application of Leibniz's Rule}

To differentiate the integral, we use Leibniz's rule:

\[
\frac{d}{dx}\left[\int_{a(x)}^{b(x)} h(x,y)\,dy\right]
=
\frac{d}{dx}b(x)\cdot h(x,b(x))
-
\frac{d}{dx}a(x)\cdot h(x,a(x))
+
\int_{a(x)}^{b(x)}\frac{\partial}{\partial x}h(x,y)\,dy.
\]

Applying this yields the convolution integral:

\[
f_Z(z) = \int_{-\infty}^{\infty} f_{XY}(z-y,y)\,dy.
\]

If $X$ and $Y$ are independent:

\[
f_Z(z) = \int_{-\infty}^{\infty} f_X(x)\,f_Y(z-x)\,dx.
\]

%------------------------------------------------

\section{Advanced Applications}

\subsection{Sum of Independent Normal Random Variables}

Assumptions:

\[
X\sim N(0,1), \qquad Y\sim N(0,1)
\]

and $X,Y$ are independent.

By substituting the Gaussian PDF into the convolution integral, it can be shown that:

\[
Z=X+Y\sim N(0,2).
\]

Thus, the sum follows a normal distribution with variance:

\[
\sigma^2 = 2.
\]

%------------------------------------------------

\subsection{Sum of Independent Exponential Random Variables}

Assumptions:  
$X$ and $Y$ are exponential random variables with parameter $\lambda$.

Method: Use convolution with

\[
f_X(x)=\lambda e^{-\lambda x},
\qquad
f_Y(y)=\lambda e^{-\lambda y}.
\]

Integral form:

\[
f_Z(z)=\int_0^z \lambda e^{-\lambda x}\cdot
\lambda e^{-\lambda(z-x)}\,dx.
\]

%------------------------------------------------

\section{Summary of Other Transformations}

\begin{itemize}
    \item Difference: $Z=X-Y$
    \item Ratio: $Z=X/Y$
    \item Square Root of Sum of Squares:
    \[
    R=\sqrt{X^2+Y^2}.
    \]
\end{itemize}

%------------------------------------------------

\end{document}
\documentclass[11pt]{article}

\usepackage{amsmath, amssymb}
\usepackage{geometry}
\geometry{margin=1in}

\title{
CSE400 -- Fundamentals of Probability in Computing\\
Lecture 11 \& 12: Transformation of Random Variables
}
\author{Instructor: Dhaval Patel, PhD}
\date{February 10, 2026}

\begin{document}

\maketitle

%------------------------------------------------

\section{Introduction and Objectives}

The primary objective of this lecture is the learning of transformation techniques for random variables.  
This involves determining the probability distribution of a new random variable $(Y)$ derived from an existing one $(X)$.

\subsection{Core Assumptions}

\begin{itemize}
    \item The Probability Density Function (PDF), $f_X(x)$, of the original random variable $X$ is known \textit{a priori}.
    \item The Cumulative Distribution Function (CDF), $F_X(x)$, of $X$ is also known.
    \item A new random variable $Y$ is defined by a functional relationship:
    \[
    y = g(x).
    \]
\end{itemize}

\subsection{Objectives}

\begin{itemize}
    \item Determine how to find the new distribution $f_Y(y)$ or $F_Y(y)$.
    \item Understand joint transformations for two random variables.
\end{itemize}

%------------------------------------------------

\section{Transformation of a Single Random Variable}

The derivation of the new PDF follows a systematic two-step process.

\subsection{Step 1: Find the Cumulative Distribution Function (CDF)}

The CDF of the new variable $Y$, $F_Y(y)$, is the probability that $Y$ is less than or equal to a value $y$:

\[
F_Y(y) = \Pr(Y \le y) = \Pr(g(x) \le y)
\]

If the function $g(x)$ is monotonically increasing:

\[
F_Y(y) = \Pr(x \le g^{-1}(y)) = F_X(g^{-1}(y))
\]

If the function $g(x)$ is decreasing:

\[
F_Y(y) = \Pr(x \ge g^{-1}(y)) = 1 - F_X(g^{-1}(y))
\]

%------------------------------------------------

\subsection{Step 2: Differentiate to Find the PDF}

To find the PDF $f_Y(y)$, differentiate the CDF with respect to $y$:

\[
f_Y(y) = \frac{d}{dy}\Big[F_X(g^{-1}(y))\Big]
\]

Using the chain rule, this results in:

\[
f_Y(y) = f_X(g^{-1}(y)) \cdot \frac{d}{dy}g^{-1}(y)
\]

Or more generally, accounting for both increasing and decreasing functions:

\[
f_Y(y) = f_X(x)\left|\frac{dx}{dy}\right|_{x=g^{-1}(y)}
\]

%------------------------------------------------

\section{Illustrative Example: Uniform to Sine Transformation}

\textbf{Problem:}  
Given $X$ is a uniformly distributed random variable on the interval $(-1,1)$, find the PDF of

\[
Y = \sin\left(\frac{\pi x}{2}\right).
\]

\subsection{Solution}

\textbf{Original PDF of $X$:}

\[
f_X(x) = \frac{1}{2}, \qquad -1 < x < 1,
\]
and $0$ otherwise.

\textbf{Inverse Function:}

\[
x = \frac{2}{\pi}\sin^{-1}(y).
\]

\textbf{Derivative:}

\[
\frac{dx}{dy} = \frac{2}{\pi}\cdot \frac{1}{\sqrt{1-y^2}}.
\]

\textbf{Resulting PDF:}

\[
f_Y(y) = \frac{1}{2}\cdot \frac{2}{\pi}\cdot \frac{1}{\sqrt{1-y^2}}
       = \frac{1}{\pi\sqrt{1-y^2}}.
\]

\textbf{Limits:}

For $x=-1$, $y=\sin(-\pi/2)=-1$.  
For $x=1$, $y=1$.  

Thus,

\[
-1 < y < 1.
\]

%------------------------------------------------

\section{Function of Two Random Variables}

This section addresses joint transformations where

\[
Z = g(X,Y).
\]

%------------------------------------------------

\subsection{Detailed Derivation: Sum of Two RVs ($Z=X+Y$)}

To find the PDF $f_Z(z)$ for the sum $Z=X+Y$:

\subsubsection{Determine the CDF}

\[
F_Z(z) = \Pr(X+Y \le z).
\]

\subsubsection{Set up the Double Integral}

\[
F_Z(z) = \int_{-\infty}^{\infty}\int_{-\infty}^{z-y}
f_{XY}(x,y)\,dx\,dy.
\]

%------------------------------------------------

\subsubsection{Application of Leibniz's Rule}

To differentiate the integral, we use Leibniz's rule:

\[
\frac{d}{dx}\left[\int_{a(x)}^{b(x)} h(x,y)\,dy\right]
=
\frac{d}{dx}b(x)\cdot h(x,b(x))
-
\frac{d}{dx}a(x)\cdot h(x,a(x))
+
\int_{a(x)}^{b(x)}\frac{\partial}{\partial x}h(x,y)\,dy.
\]

Applying this yields the convolution integral:

\[
f_Z(z) = \int_{-\infty}^{\infty} f_{XY}(z-y,y)\,dy.
\]

If $X$ and $Y$ are independent:

\[
f_Z(z) = \int_{-\infty}^{\infty} f_X(x)\,f_Y(z-x)\,dx.
\]

%------------------------------------------------

\section{Advanced Applications}

\subsection{Sum of Independent Normal Random Variables}

Assumptions:

\[
X\sim N(0,1), \qquad Y\sim N(0,1)
\]

and $X,Y$ are independent.

By substituting the Gaussian PDF into the convolution integral, it can be shown that:

\[
Z=X+Y\sim N(0,2).
\]

Thus, the sum follows a normal distribution with variance:

\[
\sigma^2 = 2.
\]

%------------------------------------------------

\subsection{Sum of Independent Exponential Random Variables}

Assumptions:  
$X$ and $Y$ are exponential random variables with parameter $\lambda$.

Method: Use convolution with

\[
f_X(x)=\lambda e^{-\lambda x},
\qquad
f_Y(y)=\lambda e^{-\lambda y}.
\]

Integral form:

\[
f_Z(z)=\int_0^z \lambda e^{-\lambda x}\cdot
\lambda e^{-\lambda(z-x)}\,dx.
\]

%------------------------------------------------

\section{Summary of Other Transformations}

\begin{itemize}
    \item Difference: $Z=X-Y$
    \item Ratio: $Z=X/Y$
    \item Square Root of Sum of Squares:
    \[
    R=\sqrt{X^2+Y^2}.
    \]
\end{itemize}

%------------------------------------------------

\end{document}
